% Français 9115, batch 13 — Gallica views 241–260.
% Pass: Opus 5 (claude-opus-5), 2026-09-19 — first pass, unchecked against
% the views by a human.
%
% What the twenty views hold. Not the Euler copy of batches 8–10: the Latin
% fair copy with its parenthesised printed page numbers has stopped somewhere
% in views 201–240 (batches 11 and 12, written in parallel with this one and
% not consulted). These leaves are a memoir of Sophie Germain's own, in
% French, in her regular fair hand with corrections: on Euler's analysis of
% the transverse vibrations of an elastic lamina held by a stylet, and on the
% case, which she says Euler did not treat, of a lamina whose two ends are
% free. Ten views carry text; the ten others are the blank versos of the same
% leaves.
%
% The leaves are alternate: an odd view is always the blank back of the leaf
% whose written side is the even view before it — view 241 is the verso of
% view 240, which belongs to batch 12, and each of views 243, 245, 247, 249,
% 251, 253, 255, 257 and 259 is the verso of the written view that precedes
% it. Nothing is written on any of them: what shows is the recto's ink read
% through the paper, reversed, including the ink number at the top left.
% None of the ten is transcribed, and none is a second exposure of another
% view: the ink numbers 118 to 127 run without a gap and no page is
% photographed twice in this batch.
%
% Each written leaf is smaller than the guard it is hinged to and is written
% on one side only. Each carries, at the head and in the middle, a short
% underlined mark — « seconde P… » at view 242, then the numbers 12, 13, 14,
% 15, 16, 17, 18, 19 and (read with doubt) 20 — which appear to be the
% articles of the memoir; and, at the top right in ink, the volume's running
% series 118 to 127.
%
%   v. 242      Head « seconde P… », the rest of the word not read. Germain
%               argues that Euler changed his mind within his own subject:
%               in the « methodus inveniendi curvas maximi minimive
%               proprietate gaudentes » he used, for the lamina free at both
%               ends, the method he kept in the memoir of 1779, yet in the
%               earlier work he rejected every solution with an odd number of
%               nodes, on the ground that the lamina would then turn about
%               the node at its middle. She sets against this Giordano
%               Riccati's « Delle Vibrazioni Sonore de cilindri » (first
%               volume of the Memorie della Società Italiana, p. 457, § XXI,
%               and p. 898 § 83), who was astonished at the opinion and found
%               that the calculation offers a way out of a conclusion
%               contrary to experiment; Riccati, she notes, seems not to have
%               known the memoir of 1779 although the volume that holds his
%               own is of 1782. Two lines at the foot are struck out whole.
%   v. 244      « Je reviens a present a 12 », marginal 7. The first of
%               Euler's solutions is general: it fixes the number of nodes by
%               the value of λ, since sin. λω = 0 makes y = 0 again at the
%               points of sin. 2λω, sin. 3λω…; the stylet may stand at any of
%               them provided n does not divide the denominator of λ. Hence
%               no regular motion, and no measurable sound, from a lamina
%               supported at both ends when the stylet is set anywhere but at
%               an aliquot division. A marginal note gives the condition that
%               would then fail. She announces the case she means to treat:
%               the stylet at an arbitrary point of a lamina free at both
%               ends.
%   v. 246      « 13 », marginal 8. « Problème » and « Solution »: the
%               equations I–IV of the memoir cited, then, for λ = 0 and
%               λ = 1, equations V–VIII; eight equations for the eight
%               coefficients α, β, γ, δ and α′, β′, γ′, δ′, and the beginning
%               of the elimination (VII × sin. ω ± VIII × cos. ω).
%   v. 248      « 14 ». Substitution into equation II, reduction by
%               cos.(1−λ)ω ∓ sin.(1−λ)ω, and the values of β′, γ′ and δ′ as
%               fractions over e^{−ω}{e^{(1−λ)ω} + cos.(1−λ)ω + sin.(1−λ)ω}.
%               Two long lines joined by a brace are set here as a
%               two-row array.
%   v. 250      « 15 ». δ′ finished; then γ and δ drawn from V and VI and
%               substituted in I, giving β, γ, δ over
%               e^{−λω} − sin. λω + cos. λω, and the first member of
%               equation III reduced to 2α{2 + (e^{λω}+e^{−λω}) cos. λω}
%               over the same denominator.
%   v. 252      « 16 ». The second member of III reduced the same way to
%               2α′[2 + (e^{(1−λ)ω}+e^{(λ−1)ω}) cos.(1−λ)ω] over its
%               denominator; equation III in its final form and the ratio
%               α/α′ that follows; then equation IV, the sum with I and II
%               divided by 2, and the resulting relation.
%   v. 254      « 17 » (the head of the leaf is torn). The comparison of the
%               two results, the transcendental equation from which λω is to
%               be drawn, and its reduction for λ = ½ into two factors. The
%               first answers to a stylet simply resting at the middle: the
%               two halves vibrate as separate laminæ of length ½a, one end
%               free and one supported, and she checks the condition against
%               cos. ω = 2e^ω/(e^{2ω}+1) to recover tang. ½ω =
%               (e^ω−1)/(e^ω+1).
%   v. 256      « 18 », marginal 10. The second factor would make the halves
%               laminæ fixed at the stylet and free at the other end, but its
%               condition agrees with the state of the ends only for ω = 0,
%               so it must be rejected — as, she says, in the case Euler
%               examined. Then the difference between the two problems: with
%               both ends supported, λ = ½ is only a particular case of a
%               general solution in which λ may be any aliquot fraction.
%   v. 258      « 19 », marginal 11. With both ends free the two conditions
%               cannot hold together except for λ = ½, and she sets out to
%               prove it: the two quantities are transformed, the value drawn
%               from the second is put into the first, and the whole is
%               multiplied by (e^{λω}−e^{−λω})(e^{(1−λ)ω}−e^{(λ−1)ω}). A last
%               line is struck out whole.
%   v. 260      « 20 » (read with doubt), marginal 12. The equation reduced
%               to 2 sin. λω {(e^{(1−2λ)ω}+e^{(2λ−1)ω}) sin. ω −
%               (e^ω−e^{−ω}) cos. ω} = 0; sin. λω cannot vanish, so the other
%               factor must, which gives a value of tang. ω the state of the
%               ends contradicts unless λ = ½. She adds that the same holds
%               when the ends are fixed instead of free, and suppresses those
%               calculations « a cause de leur parfaite analogie avec les
%               précédans ». The leaf ends on « a l'équation. » and the piece
%               runs on past this batch.
%
% The numbers on the leaves, and why no \folio{} is written. Two kinds.
% (1) The volume's running ink series, at the top right of every written
% leaf: 118 at view 242, 119 at 244, 120 at 246, 121 at 248, 122 at 250,
% 123 at 252, 124 at 254, 125 at 256, 126 at 258, 127 at 260 — one per leaf,
% without a gap, and the continuation of the series that stood at 97 on view
% 199 in batch 10. It is ink, not the Bibliothèque's pencil, and the earlier
% batches of this volume wrote no \folio{} for it; this one does not either.
% On the blank versos the same figure shows through from the recto, reversed,
% at the top left; it is not written there. (2) The memoir's own article
% numbers, underlined, at the head and in the middle of the written leaf:
% « seconde P… » at view 242, then 12, 13, 14, 15, 16, 17, 18, 19, 20 at
% views 244 to 260. At view 254 the top of the leaf is torn and only the foot
% of the figure survives; at view 260 the second figure is shaped like an ω.
% Both are marked \uncertain{}. A third series, of small figures alone in the
% left margin — 7 (v. 244), 8 (v. 246), 10 (v. 256), 11 (v. 258), 12 (v. 260)
% — numbers the paragraphs; 9 is missing, and the struck mark that opens
% view 256 may be it. Above the torn top edge of view 260, on the mount and
% not on the leaf, an underlined 21 is written. Every number above was read
% on its view; none is inferred, and no pencilled foliation of the
% Bibliothèque was seen anywhere in the batch.
%
% Notation. Germain's letters follow the Euler copy of batches 8–10:
% α, β, γ, δ and α′, β′, γ′, δ′ for the eight constants, ω for the angle,
% λ for the fraction of the length at which the stylet stands (a small hooked
% letter crossed by a stroke, as in the copy), a for the whole length.
% « sin. », « cos. » and « tang. » are written with a stop, which is kept;
% the stop is set where it was seen. Exponentials are written e with the
% exponent above and the sign of the exponent often below the line; they are
% set as ordinary exponents. Braces { } and brackets [ ] are the writer's.
% Her spelling stays: « primitivee », « dela », « ceque », « amoins »,
% « aulieu », « cidessus », « c'est adire », « desorte », « appuié »,
% « suprime », « accolade »-less run-on sentences and all. Where the leaf and
% the calculation disagree the leaf stands, with a note: view 260 writes
% « un nombre impair » twice where the second should be « pair ».
%
% What could not be read. Twelve places are marked \ill{}, and all but two
% are under a blot or a deletion. View 242 has six: the end of the heading;
% the struck word opening the text, of which only « Qu » is formed; the last
% word of the Riccati line, after « de »; two short stretches inside the long
% struck passage; and one or two words under a blot in the struck sentence at
% the foot. View 246 has one, a struck word or two before « comme dans le
% cas ». View 256 has four: the mark, perhaps a figure, struck at the head of
% the first line, and three blots struck in the last lines. View 258 has one:
% its last line, struck out from end to end.
%
% Readings offered with doubt are marked \uncertain{}: « ait »,
% « Giordanono » at view 242; « il », « & » at view 244; « de même » at
% view 246; the figure 4 at views 250 and 252; « 17 » and the reference
% « V. cas. 1er » at view 254; « Au » and « avant » and the « cos. ω » under
% a blot at view 258; « 20 » and « concluroit » at view 260. Nothing was
% guessed to fill a gap.
%
% Nothing here was checked against any published transcription or edition of
% Germain, of Euler or of Riccati.
\documentclass[11pt,a4paper]{article}
% The preamble shared by every transcription of the mathematicians' manuscripts in Gallica.
%
% It rests on one idea: **uncertainty compounds, it does not resolve.** A
% transcription that smooths over a doubtful word has destroyed the only thing
% it was made for — a reader must be able to tell, at every point, what was on
% the page and what was supplied. The apparatus macros below are therefore the
% critical apparatus, not decoration.
%
% The same commands are recognised by `scripts/render.mjs`, which produces the
% site's reading view, and by `scripts/tei.mjs`. A macro added here must be
% added there too, or rendering fails loudly — which is the wanted behaviour.
%
% Comments here are English, like the rest of the repository. The documents
% this preamble sets are French, and that is deliberate: see the skills.

\ProvidesPackage{germain}[2026/09/15 Transcriptions of mathematicians' manuscripts in Gallica]

\usepackage{amsmath,amssymb,amsthm}
\usepackage{mathrsfs}
% Kept from the parent preamble so the renderer's diagram support stays
% available; these manuscripts draw no commutative diagrams, and nothing here uses it.
\usepackage{tikz-cd}
\usepackage[english,french]{babel}
\usepackage{fontspec}

% --- Greek ------------------------------------------------------------------
% The text face stays fontspec's default, Latin Modern, which has no Greek:
% XeTeX drops every glyph it lacks with a line in the log and nothing on the
% page, and the Greek words the manuscripts quote vanished from the PDFs. So
% the Greek and Greek Extended blocks get a XeTeX character class of their own,
% and entering or leaving a run of it switches the family to CMU Serif — the
% same Computer Modern design, with polytonic Greek — and back. Only the
% family changes: a Greek word in \textit or \textbf keeps its shape and series.
% The transitions to and from every other class carry the tokens that class
% has towards an ordinary letter, so French spacing before « ; » or inside
% « » still holds after a Greek word. No grouping, so no brace or \endgroup
% can come out unbalanced; maths is left alone, since XeTeX inserts nothing
% there. Kept identical in `exercices.sty`.
\makeatletter
\newfontfamily\germainGreekfont{cmun}[
  Extension=.otf, NFSSFamily=germaingreek,
  UprightFont=*rm, ItalicFont=*ti, SlantedFont=*sl,
  BoldFont=*bx, BoldItalicFont=*bi, BoldSlantedFont=*bl]
\newXeTeXintercharclass\germain@greek
\def\germain@greekrange#1#2{%
  \@tempcnta=#1\relax
  \loop
    \XeTeXcharclass\@tempcnta=\germain@greek
    \ifnum\@tempcnta<#2\advance\@tempcnta\@ne\repeat}
\germain@greekrange{"0370}{"03FF}
\germain@greekrange{"1F00}{"1FFF}
\def\germain@greekprev{\rmdefault}
\def\germain@greekin{%
  \edef\germain@greekprev{\f@family}\fontfamily{germaingreek}\selectfont}
\def\germain@greekout{\fontfamily{\germain@greekprev}\selectfont}
\def\germain@greekjoin{%
  \XeTeXinterchartokenstate=\@ne
  \@tempcnta=\z@
  \loop
    \ifnum\@tempcnta=\germain@greek\else
      \XeTeXinterchartoks\germain@greek\@tempcnta=\expandafter{%
        \expandafter\germain@greekout\the\XeTeXinterchartoks\z@\@tempcnta}%
      \XeTeXinterchartoks\@tempcnta\germain@greek=\expandafter{%
        \the\XeTeXinterchartoks\@tempcnta\z@\germain@greekin}%
    \fi
    \ifnum\@tempcnta<4095\advance\@tempcnta\@ne\repeat}
% babel-french sets its own classes, and saves and restores the interchar
% state, each time a language is selected: join again after it does.
\addto\extrasfrench{\germain@greekjoin}
\addto\extrasenglish{\germain@greekjoin}
\AtBeginDocument{\germain@greekjoin}
\makeatother

\usepackage{geometry}
\geometry{a4paper,margin=25mm,marginparwidth=28mm}
\usepackage{marginnote}
\usepackage{xcolor}
\usepackage[normalem]{ulem}
\setlength{\emergencystretch}{3em}
\usepackage{hyperref}
\hypersetup{colorlinks=true,linkcolor=black,urlcolor=black,pdfborder={0 0 0}}

% --- Batch metadata ---------------------------------------------------------
% Taken as they stand by the rendering script, which shows them at the head of
% the reading view. They fix what the file claims to cover. The macro names are
% the parent project's, so that the scripts stay shared; \folder here means the
% volume's slug — `fr-9115` — and \pages counts Gallica views.

\newcommand{\folder}[1]{\gdef\grothFolder{#1}}
\newcommand{\batch}[1]{\gdef\grothBatch{#1}}
\newcommand{\pages}[2]{\gdef\grothFirst{#1}\gdef\grothLast{#2}}
\newcommand{\foldertitle}[1]{\gdef\grothFoldertitle{#1}}

% The holder's own shelfmark, printed as they write it: « Français 9115 ».
\newcommand{\shelfmark}[1]{\gdef\grothShelfmark{#1}}

% The dating, copied verbatim from the holder's catalogue — never inferred
% here. « XIXe siècle » is the BnF's; a pass that dates a leaf from its content
% says so in a \note{}, as its own inference.
\newcommand{\dating}[1]{\gdef\grothDating{#1}}

% The status notice, printed in the title block and repeated as a diagonal
% watermark on every page of the PDF. The manuscripts are in the public
% domain; what this line declares is not a right but a quality — an
% unverified first pass by a machine. The skills write the line; a file
% without it gets no watermark, so leaving it out is visible in the output.
\newcommand{\watermark}[1]{\gdef\grothWatermark{#1}}

\gdef\grothFolder{?}\gdef\grothBatch{}\gdef\grothFirst{?}\gdef\grothLast{?}%
\gdef\grothFoldertitle{}\gdef\grothShelfmark{}\gdef\grothDating{}\gdef\grothWatermark{}

% --- The résumé -------------------------------------------------------------
\newenvironment{resume}
  {\small\setlength{\parskip}{0.45em}}
  {\par\medskip\hrule\medskip}

% --- Keywords ---------------------------------------------------------------
% In English, closing the modernised reading's résumé: the modern vocabulary
% under which to search for what the leaves construct. The manifest extracts
% them to tag the volume — this line is the single source of the tags.
\newcommand{\keywords}[1]{\par\smallskip\noindent
  {\footnotesize\textsc{Keywords} --- \textit{#1}}\par}

% --- The critical apparatus -------------------------------------------------

% The Gallica view number, in the margin. This is the anchor that lets a
% disagreement over a formula be settled by looking at *one* image rather than
% a volume of 750. The site uses it to turn the facsimile as the transcription
% is scrolled. A view is one scanned image: usually one side of a leaf.
\newcommand{\page}[1]{%
  \marginnote{\footnotesize\color{gray!70}\textsf{v.\,#1}}%
  \label{page:#1}\ignorespaces}

% The BnF's pencilled foliation, where the leaf carries it and a pass can read
% it: « 198r ». This is what the literature cites — « ff. 198r–208v » — and
% the one line that turns a citation into a view. Recorded, never inferred:
% a folio a pass cannot read is not written.
\newcommand{\folio}[1]{%
  \marginnote{\footnotesize\color{gray!50}\textsf{f.\,#1}}\ignorespaces}

% The modernised reading's coarser counterpart to \page{}: which views a
% section is drawn from.
\newcommand{\pagerange}[2]{%
  \marginnote{\footnotesize\color{gray!70}\textsf{v.\,#1--#2}}%
  \ignorespaces}

% Illegible: never guessed.
\newcommand{\ill}{\textcolor{red!60!black}{[\,\ldots\,]}}

% A doubtful reading: the word is given, and flagged as uncertain.
\newcommand{\uncertain}[1]{\uline{#1}}

% An editorial addition — a missing letter, an implied word.
\newcommand{\add}[1]{\textcolor{blue!50!black}{[#1]}}

% Struck out by the writer. What was crossed out is part of the document.
\newcommand{\struck}[1]{\sout{#1}}

% The transcriber's note, on the state of the page or the mathematics.
\newcommand{\note}[1]{\par\smallskip{\small\color{gray!60!black}\textit{[#1]}}\par\smallskip}

% A marginal note of hers, distinct from the body of the text.
\newcommand{\marginal}[1]{\par\smallskip\noindent\hspace{1em}%
  {\small\color{gray!60!black}#1}\par\smallskip}

% --- Title ------------------------------------------------------------------
% The title block is French, like the documents themselves.

\AddToHook{shipout/background}{%
  \ifx\grothWatermark\empty\else
    \begin{tikzpicture}[remember picture, overlay]
      \node[rotate=52, scale=3.4, align=center, text=gray!18,
            font=\sffamily\bfseries]
        at (current page.center) {\grothWatermark};
    \end{tikzpicture}%
  \fi}

\AtBeginDocument{%
  \begin{center}
    {\large\bfseries Manuscrits de mathématiciens (Gallica) ---
      \ifx\grothShelfmark\empty\grothFolder\else\grothShelfmark\fi}\\[2pt]
    {\itshape\grothFoldertitle}\\[4pt]
    {\small vues \grothFirst--\grothLast
      \ifx\grothBatch\empty\else\ (lot \grothBatch)\fi}\\[2pt]
    \ifx\grothDating\empty\else
      {\small\color{gray!60!black}Datation du catalogue : \grothDating}\\[2pt]
    \fi
    \ifx\grothWatermark\empty\else
      {\small\bfseries\color{red!45!gray}\grothWatermark}\\[2pt]
    \fi
    {\footnotesize\color{gray!60!black}Transcription établie d'après les images IIIF de Gallica.
     Source gallica.bnf.fr / Bibliothèque nationale de France.\\
     Les lectures incertaines sont soulignées, les illisibles notées [\,\ldots\,],
     les ajouts entre crochets ; \textsf{v.} numérote les vues de Gallica,
     \textsf{f.} la foliotation au crayon de la BnF.}
  \end{center}
  \vspace{1em}}

\endinput


\folder{fr-9115}
\shelfmark{Français 9115}
\batch{13}
\pages{241}{260}
\dating{1801-1900}
\watermark{Lecture automatique\\non vérifiée}
\foldertitle{Papiers de Sophie GERMAIN. — Recueil de dissertations et problèmes mathématiques et physiques.}

\begin{document}

\note{Numérotation des feuillets. Les feuillets de ce lot sont écrits d'un
seul côté : les vues 241, 243, 245, 247, 249, 251, 253, 255, 257 et 259 sont
les versos blancs des feuillets dont le recto est la vue paire qui précède
(la vue 241 est le verso de la vue 240, qui appartient au lot 12) ; on n'y
voit que l'encre du recto par transparence, renversée, et elles ne sont pas
transcrites. Chaque feuillet écrit porte en haut à droite, à l'encre, un
nombre de la série courante du volume : 118 (vue 242), 119 (244), 120 (246),
121 (248), 122 (250), 123 (252), 124 (254), 125 (256), 126 (258), 127 (260),
sans lacune ; c'est la suite de la série qui était à 97 à la vue 199. Elle est
à l'encre et non au crayon de la Bibliothèque, et aucune foliotation n'a été
lue sur ces vues : aucun « f. » n'est donc porté ici. Un second nombre,
souligné, est écrit en tête et au milieu de chaque feuillet — « seconde P… »
(vue 242), puis 12, 13, 14, 15, 16, 17, 18, 19, 20 (vues 244 à 260) : ce sont
les articles du mémoire. Un troisième, de petits chiffres seuls dans la marge
de gauche, numérote les alinéas : 7 (244), 8 (246), 10 (256), 11 (258), 12
(260). Tous ces nombres ont été lus sur la vue ; aucun n'est déduit.}

\page{242}
seconde P\ill{}

D'ailleurs

\struck{\ill{}} quelque repugnance que l'on puisse avoir a devier de l'opinion
d'Euler on est pourtant forcé de convenir que sans sortir du sujet de ce
memoire il a lui même varié de sentiment, en effet quoique dans sa
\emph{methodus inveniendi curvas maximi minimive proprietate gaudentes}, il
\uncertain{ait} suivi, pour l'analyse du cas ou les deux extremités dela lame
elastique vibrante sont libres, la même methode qu'il a egalement adopté dans
son memoire que j'ai déja cité une methode qui ne differe pas essentiellement
de celle qu'il a depuis adopté dans \struck{le memoire que j'ai déja fait} son
memoire de 1779, il s'est pourtant decidé \struck{a rejetter} dans son premier
travail a rejetter toutes les solutions qui donnent lieu a un nombre impair de
nœuds la raison qu'il en aporte est qu'il doit exister alors un mouvement
giratoire de la lame autour du nœud placé au milieu de cette \add{même} lame.
Giordano Ricati P. 898 §. 83 de \ill{}

M\textsuperscript{r} \uncertain{Giordanono} Ricati dans son memoire
\emph{Delle Vibrazioni Sonore de cilindri} ouvrage dans lequel il suit
d'ailleurs les principes qu'Euler a etabli dans \struck{son \add{premier}
traité} son livre de \emph{methodus} \&c. \struck{et dans lequel il se trouve
nommément le cas des extremités libres \ill{} dela courbe elastique.}
\struck{l'ouvrage que je viens de citer} s'étonne de cette opinion
\struck{les travaux d'Euler et il calcule avec plus de precision les valeurs
qui \ill{}} et remarque que le calcul fournit des ressources pour échaper a une
conclusion qui est \struck{aussi} contraire a l'experience. V. le premier
volume des memoires de la société italienne \struck{P. 457} §. \struck{XXI}
P. 457. M\textsuperscript{r} Ricati ne paroit pas avoir eu connoissance du
memoire 1779 quoique le volume qui renferme le sien soit de 1782 il \add{y}
auroit vû qu'Euler n'admet l'existence des solutions qui donnent lieu a un
nombre impair de nœud et que \struck{dans le scholie} le seul usage qu'il
fasse a ce sujet de l'exclusion du mouvement giratoire \struck{le réduit}
\struck{qu'}a prouver que dans le mouvement le plus simple dela lame elastique
dont les deux extremités sont libres la courbe doit couper l'axe en deux
points. V. le Scholie §. 81

\struck{J'ai rapporté ces exemples afin d'écarter en quelque sorte les doutes
que m'a \ill{} inspiré les parties des memoires qui font l'objet de ces
remarques}
\note{Le feuillet, plus petit que la garde, est monté sur un onglet ; il n'est
écrit que sur la moitié supérieure. En tête, au milieu et souligné : « seconde »
suivi d'un P majuscule et de deux traits épais repassés, qui n'ont pas été lus ;
le sens appellerait « Partie », mais rien n'est déduit ici. En haut à droite, à
l'encre, le nombre 118. Le mot biffé qui ouvre le texte commence par « Qu » ;
la suite n'est pas lue. Dans « au milieu de cette même lame », « même » est
ajouté au-dessus de la ligne avec un signe d'insertion. Un long trait
horizontal traverse « Giordano Ricati » : il part de la barre du t de
« Ricati » et pourrait n'être que cette barre rejetée en arrière, comme
ailleurs dans cette main ; il n'est donc pas donné pour une biffure. Le dernier
mot de cette ligne, après « de », n'est pas lu, et la phrase est reprise à la
ligne suivante. Entre « qu'Euler » et « n'admet », un petit signe en croix,
sans addition au-dessus. Les deux dernières lignes du feuillet sont biffées
d'un bout à l'autre ; une tache d'encre y couvre un ou deux mots après « que
m'a ». Les titres latins et italiens sont soulignés dans le manuscrit.}

\page{244}
\marginal{7}
Je reviens a present a \emph{12}

A l'egard de la premiere solution \add{d'Euler} on voit qu'elle est générale
elle montre comment il s'etablit un nombre de nœuds dependans dela valeur de
$\lambda$ car puisque \add{on a} $y=0$ par la supposition
$\text{sin.}\,\lambda\omega=0$ on aura encore $y=0$ pour les points qui
correspondent au sinus \struck{sinus} multiples, tels que
$\text{sin.}\,2\lambda\omega$, $\text{sin.}\,3\lambda\omega$ \&c.

D'ou il resulte que le mouvement dela lame est le même soit que le stilet soit
placé a l'un ou l'autre de ces points pourvu toutefois qu'en designant par
$\text{sin.}\,n\lambda\omega$ un des sinus de cette suite $n$ ne soit pas
diviseur du denominateur dela fraction $\lambda$, car \uncertain{il} resulte
encore de cette solution que tout mouvement regulier et par consequent tout
son mesurable est impossible a obtenir dela lame dont les deux extremités sont
appuiées lorsque le stilet est placé a tout autre point qu'a une division
aliquote dela même lame, on voit en effet qu'excepté ce cas les conditions
$\text{sin.}\,\omega=0$ et $\text{sin.}\,\lambda\omega=0$ n'ont plus lieu en
mêmes tems et que l'equation du problème n'est plus
$\text{sin.}\,\lambda\omega$ pour facteur.
\marginal{alors la supposition $\text{sin.}\,n\lambda\omega=0$ n'entraineroit
pas necessairement $\text{sin.}\,\lambda\omega=0$}

Ces remarques etablissent le plus grand rapport entre les circonstances du
mouvement d'une telle lame et celles qui ont lieu pour la corde vibrante,
malheureusement ce cas des deux extremités appuiées dont l'analyse est le plus
facile n'est pas celui qui se presente le plus commodement dans l'experience
\uncertain{\&} la pratique.

Je me propose d'appliquer \add{en suite} la méthode d'Euler aux cas ou le
stilet est placé dans un point quelconque de l'etendue d'une lame elastique
vibrante dont les deux extremités sont libres, l'importance de ce cas m'engage
a developper \add{ici} les calculs du problème qui y est relatif.
\note{En tête : le chiffre 7 seul dans la marge de gauche ; puis, d'une écriture
plus fine, « Je reviens a present a », et au milieu de la page le nombre 12
souligné, qui paraît être le terme de cette phrase. En haut à droite, à
l'encre, le nombre 119. Le mot grec noté $\lambda$ est une petite lettre
crochue barrée d'un trait, comme dans la copie latine des lots précédents.
« d'Euler » et « on a » sont ajoutés au-dessus de la ligne. Entre « car » et
« resulte » se dresse un grand signe à deux jambages qui relie la ligne à la
note marginale ; il porte peut-être le « il » qui manque à la phrase. Dans
« n'est plus », « est » est repassé sur un mot plus court. Entre
« l'experience » et « la pratique » un signe court empâté, lu avec doute
« \& ». Au-dessus de « d'appliquer », avant « en suite », un mot biffé que
l'on ne lit pas avec certitude, peut-être « egalement ». Le feuillet, plus
petit que la garde, est monté sur un onglet et s'arrête sur « relatif ».}

\page{246}
\emph{13}

\section{Problème}

\marginal{8}
Si une lame elastique \struck{est} libre a l'une et l'autre de ses extremités
est soumise a l'action d'un stilet dans un point quelconque de son etendue;
trouver tous les mouvemens de vibration dont elle est susceptible?

\subsection{Solution}

Je reprends les équations

\[ \text{I.}\quad \alpha e^{\lambda\omega}+\beta e^{-\lambda\omega}+\gamma\,\text{sin.}\,\lambda\omega+\delta\cos.\lambda\omega=0 \]
\[ \text{II.}\quad \alpha' e^{\lambda\omega}+\beta' e^{-\lambda\omega}+\gamma'\,\text{sin.}\,\lambda\omega+\delta'\cos.\lambda\omega=0 \]
\[ \text{III.}\quad \alpha e^{\lambda\omega}-\beta e^{-\lambda\omega}+\gamma\cos.\lambda\omega-\delta\,\text{sin.}\,\lambda\omega=\alpha' e^{\lambda\omega}-\beta' e^{-\lambda\omega}+\gamma'\cos.\lambda\omega-\delta'\,\text{sin.}\,\lambda\omega \]
\[ \text{IV.}\quad (\alpha-\alpha')e^{\lambda\omega}+(\beta-\beta')e^{-\lambda\omega}-(\gamma-\gamma')\,\text{sin.}\,\lambda\omega-(\delta-\delta')\cos.\lambda\omega=0 \]

Donnés par Euler dans l'analyse du même problème relatif au cas ou les
extremités de la lame sont tout l'une et l'autre appuiées par ce qu'il est
clair que ces équations sont egalement applicables au cas present puisqu'elles
sont independantes de l'état des extremités.

Maintenant donc, en considerant l'état des extremités (Voyez cas 1\textsuperscript{er}
N\textsuperscript{o} 5 du memoire cité) il en resultera quatre équations) et
faisant successivement $\lambda=0$ et $\lambda=1$, savoir

\[ \text{V.}\quad \alpha-\beta=\gamma, \qquad \text{VI.}\quad \alpha+\beta=\delta \]
\[ \text{VII.}\quad \alpha' e^{\omega}+\beta' e^{-\omega}-\gamma'\,\text{sin.}\,\omega-\delta'\cos.\omega=0 \quad\text{et}\quad \text{VIII.}\quad \alpha' e^{\omega}-\beta' e^{-\omega}-\gamma'\cos.\omega+\delta'\,\text{sin.}\,\omega=0 \]

ainsi on \struck{\ill{}} comme dans le cas des extremités appuiées, huit
equations pour déterminer les huit coëfficiens $\alpha$, $\beta$, $\gamma$,
$\delta$ et $\alpha'$, $\beta'$, $\gamma'$, $\delta'$. et l'on parvient
\uncertain{de même} a une équation finale delaquelle on peut tirer la valeur
d'$\omega$.

Pour cela je commence par multiplier l'equation VII par $\text{sin.}\,\omega$
et l'equation VIII par $\cos.\omega$ puis je les ajoute l'une a l'autre ce qui
donne
\[ \alpha' e^{\omega}(\text{sin.}\,\omega+\cos.\omega)+\beta' e^{-\omega}(\text{sin.}\,\omega-\cos.\omega)=\gamma' \]
je multiplie en suite l'equation VII par $\cos.\omega$ et l'equation VIII par
$\text{sin.}\,\omega$ ce qui donne en les retranchant l'une de l'autre
\[ \alpha' e^{\omega}(\cos.\omega-\text{sin.}\,\omega)+\beta' e^{-\omega}(\text{sin.}\,\omega+\cos.\omega)=\delta' \]

En substituant les presentes valeurs de $\gamma'$ et $\delta'$ dans l'equation
\note{En tête, au milieu et souligné, le nombre 13 ; en haut à droite, à
l'encre, le nombre 120. Le chiffre 8 est seul dans la marge de gauche, devant
« Si » ; un signe court, peut-être un chiffre, le suit. « Problème » et
« Solution » sont soulignés. Dans « ainsi on … comme dans le cas », un ou deux
mots sont biffés et repassés et ne sont pas lus ; le sens demanderait « on a ».
« huit » et, plus haut, « tous » sont repassés à la plume. Dans « parvient de
même », un pâté couvre le début de « de même ». Le $\gamma$ isolé des
équations V et VII est tracé en boucle renversée, différente du $\gamma$ qui
précède « sin. » dans les équations I à IV. Le feuillet, plus petit que la
garde, est monté sur un onglet ; une tache d'encre est au bas de la marge.
Le texte s'arrête sur « dans l'equation » ; la suite est à la vue 248.}

\page{248}
\emph{14}

Dans l'équation II elle devient

\[ \text{II}\quad \alpha' e^{\lambda\omega}+\beta' e^{-\lambda\omega}+\{\alpha' e^{\omega}(\text{sin.}\,\omega+\cos.\omega)+\beta' e^{-\omega}(\text{sin.}\,\omega-\cos.\omega)\}\,\text{sin.}\,\lambda\omega \]
\[ +\{\alpha' e^{\omega}(\cos.\omega-\text{sin.}\,\omega)+\beta' e^{-\omega}(\text{sin.}\,\omega+\cos.\omega)\}\cos.\lambda\omega=0 \]

ou
\[ \alpha'\{e^{\lambda\omega}+e^{\omega}(\text{sin.}\,\omega\,\text{sin.}\,\lambda\omega+\text{sin.}\,\lambda\omega\cos.\omega+\cos.\lambda\omega\cos.\omega-\text{sin.}\,\omega\cos.\lambda\omega)\} \]
\[ +\beta'\{e^{-\lambda\omega}+e^{-\omega}(\text{sin.}\,\omega\,\text{sin.}\,\lambda\omega-\cos.\omega\,\text{sin.}\,\lambda\omega+\text{sin.}\,\omega\cos.\lambda\omega+\cos.\lambda\omega\cos.\omega)\}=0 \]

ou enfin
\[ \alpha'\{e^{\lambda\omega}+e^{\omega}(\cos.(1-\lambda)\omega-\text{sin.}\,(1-\lambda)\omega)\}+\beta'\{e^{-\lambda\omega}+e^{-\omega}(\cos.(1-\lambda)\omega+\text{sin.}\,(1-\lambda)\omega)\}=0 \]

D'ou on tire
\[ \beta'=-\frac{\alpha' e^{\omega}\{e^{(\lambda-1)\omega}+\cos.(1-\lambda)\omega-\text{sin.}\,(1-\lambda)\omega\}}{e^{-\omega}\{e^{(1-\lambda)\omega}+\cos.(1-\lambda)\omega+\text{sin.}\,(1-\lambda)\omega\}} \]

par consequent

\[ \gamma'=\alpha' e^{\omega}(\text{sin.}\,\omega+\cos.\omega)+\beta' e^{-\omega}(\text{sin.}\,\omega-\cos.\omega) \]
\[ =\frac{\alpha'(\text{sin.}\,\omega+\cos.\omega)\{e^{(1-\lambda)\omega}+\cos.(1-\lambda)\omega+\text{sin.}\,(1-\lambda)\omega\}}{e^{-\omega}\{e^{(1-\lambda)\omega}+\cos.(1-\lambda)\omega+\text{sin.}\,(1-\lambda)\omega\}} \]
\[ -\frac{\alpha'(\text{sin.}\,\omega-\cos.\omega)\{e^{(\lambda-1)\omega}+\cos.(1-\lambda)\omega-\text{sin.}\,(1-\lambda)\omega\}}{e^{-\omega}\{e^{(1-\lambda)\omega}+\cos.(1-\lambda)\omega+\text{sin.}\,(1-\lambda)\omega\}} \]
\[ =\frac{\alpha'\{(e^{(1-\lambda)\omega}-e^{(\lambda-1)\omega})\,\text{sin.}\,\omega+(e^{(1-\lambda)\omega}+e^{(\lambda-1)\omega})\cos.\omega\}}{e^{-\omega}\{e^{(1-\lambda)\omega}+\cos.(1-\lambda)\omega+\text{sin.}\,(1-\lambda)\omega\}} \]
\[ +\frac{\left\{\begin{array}{l} \text{sin.}\,\omega\cos.(1-\lambda)\omega+\text{sin.}\,\omega\,\text{sin.}\,(1-\lambda)\omega+\cos.\omega\cos.(1-\lambda)\omega+\cos.\omega\,\text{sin.}\,(1-\lambda)\omega \\ -\text{sin.}\,\omega\cos.(1-\lambda)\omega+\text{sin.}\,\omega\,\text{sin.}\,(1-\lambda)\omega+\cos.\omega\cos.(1-\lambda)\omega-\cos.\omega\,\text{sin.}\,(1-\lambda)\omega \end{array}\right\}\alpha'}{e^{-\omega}\{e^{(1-\lambda)\omega}+\cos.(1-\lambda)\omega+\text{sin.}\,(1-\lambda)\omega\}} \]

ou enfin
\[ \gamma'=\frac{\alpha'\{(e^{(1-\lambda)\omega}-e^{(\lambda-1)\omega})\,\text{sin.}\,\omega+(e^{(1-\lambda)\omega}+e^{(\lambda-1)\omega})\cos.\omega+2\cos.\lambda\omega\}}{e^{-\omega}\{e^{(1-\lambda)\omega}+\cos.(1-\lambda)\omega+\text{sin.}\,(1-\lambda)\omega\}} \]

on trouve de même

\[ \delta'=\alpha' e^{\omega}(\cos.\omega-\text{sin.}\,\omega)+\beta' e^{-\omega}(\text{sin.}\,\omega+\cos.\omega) \]
\[ =\frac{\alpha'(\cos.\omega-\text{sin.}\,\omega)\{e^{(1-\lambda)\omega}+\cos.(1-\lambda)\omega+\text{sin.}\,(1-\lambda)\omega\}}{e^{-\omega}\{e^{(1-\lambda)\omega}+\cos.(1-\lambda)\omega+\text{sin.}\,(1-\lambda)\omega\}} \]
\[ -\frac{\alpha'(\text{sin.}\,\omega+\cos.\omega)\{e^{(\lambda-1)\omega}+\cos.(1-\lambda)\omega-\text{sin.}\,(1-\lambda)\omega\}}{e^{-\omega}\{e^{(1-\lambda)\omega}+\cos.(1-\lambda)\omega+\text{sin.}\,(1-\lambda)\omega\}} \]
\[ =\frac{\alpha'\{(e^{(1-\lambda)\omega}-e^{(\lambda-1)\omega})\cos.\omega-(e^{(1-\lambda)\omega}+e^{(\lambda-1)\omega})\,\text{sin.}\,\omega\}}{e^{-\omega}\{e^{(1-\lambda)\omega}+\cos.(1-\lambda)\omega+\text{sin.}\,(1-\lambda)\omega\}} \]
\[ +\frac{\alpha'\{(\cos.\omega-\text{sin.}\,\omega)(\cos.(1-\lambda)\omega+\text{sin.}\,(1-\lambda)\omega)-(\text{sin.}\,\omega+\cos.\omega)(\cos.(1-\lambda)\omega-\text{sin.}\,(1-\lambda)\omega)\}}{e^{-\omega}\{e^{(1-\lambda)\omega}+\cos.(1-\lambda)\omega+\text{sin.}\,(1-\lambda)\omega\}} \]
\note{En tête, au milieu et souligné, le nombre 14 ; en haut à droite, à l'encre,
le nombre 121. Le feuillet ne porte que des calculs, sans autre prose que les
liaisons « ou », « ou enfin », « par consequent », « on trouve de même ». Dans
« ou enfin » qui précède la valeur définitive de $\gamma'$, le mot « ou » est
traversé d'un trait. Les deux longues lignes réunies par une accolade sont
écrites l'une sous l'autre, l'accolade ouvrante et l'accolade fermante
embrassant les deux, et le facteur $\alpha'$ est porté après l'accolade
fermante ; elles sont ici rendues par un tableau à deux lignes. Le feuillet,
plus petit que la garde, est monté sur un onglet ; il s'arrête sur cette
dernière fraction, et le calcul se poursuit à la vue 250.}

\page{250}
\emph{15}

enfin
\[ \delta'=\frac{\alpha'\{(e^{(1-\lambda)\omega}-e^{(\lambda-1)\omega})\cos.\omega-(e^{(1-\lambda)\omega}+e^{(\lambda-1)\omega})\,\text{sin.}\,\omega-2\,\text{sin.}\,\lambda\omega\}}{e^{-\omega}\{e^{(1-\lambda)\omega}+\cos.(1-\lambda)\omega+\text{sin.}\,(1-\lambda)\omega\}} \]

Je prends ensuite les valeurs de $\gamma$ et $\delta$ données par les équations
V et VI pour les substituer dans l'equation

\[ \text{I.}\quad \alpha e^{\lambda\omega}+\beta e^{-\lambda\omega}+\gamma\,\text{sin.}\,\lambda\omega+\delta\cos.\lambda\omega=0 \]

elle devient par là

\[ \alpha e^{\lambda\omega}+\beta e^{-\lambda\omega}+(\alpha-\beta)\,\text{sin.}\,\lambda\omega+(\alpha+\beta)\cos.\lambda\omega=0 \]

ou
\[ \alpha\{e^{\lambda\omega}+\text{sin.}\,\lambda\omega+\cos.\lambda\omega\}+\beta\{e^{-\lambda\omega}-\text{sin.}\,\lambda\omega+\cos.\lambda\omega\}=0 \]

D'ou on tire
\[ \beta=-\frac{\alpha\{e^{\lambda\omega}+\text{sin.}\,\lambda\omega+\cos.\lambda\omega\}}{e^{-\lambda\omega}-\text{sin.}\,\lambda\omega+\cos.\lambda\omega} \]

et par consequent

\[ \gamma=\alpha-\beta \]
\[ =\frac{\alpha\{e^{-\lambda\omega}-\text{sin.}\,\lambda\omega+\cos.\lambda\omega+e^{\lambda\omega}+\text{sin.}\,\lambda\omega+\cos.\lambda\omega\}}{e^{-\lambda\omega}-\text{sin.}\,\lambda\omega+\cos.\lambda\omega} \]
\[ =\frac{\alpha\{e^{\lambda\omega}+e^{-\lambda\omega}+2\cos.\lambda\omega\}}{e^{-\lambda\omega}-\text{sin.}\,\lambda\omega+\cos.\lambda\omega} \]

\[ \delta=\alpha+\beta \]
\[ =\frac{\alpha\{e^{-\lambda\omega}-\text{sin.}\,\lambda\omega+\cos.\lambda\omega-e^{\lambda\omega}-\text{sin.}\,\lambda\omega-\cos.\lambda\omega\}}{e^{-\lambda\omega}-\text{sin.}\,\lambda\omega+\cos.\lambda\omega} \]
\[ =-\frac{\alpha\{e^{\lambda\omega}-e^{-\lambda\omega}+2\,\text{sin.}\,\lambda\omega\}}{e^{-\lambda\omega}-\text{sin.}\,\lambda\omega+\cos.\lambda\omega} \]

Il faut a present substituer les valeurs que l'on vient de trouver pour
$\alpha$, $\beta$, $\gamma$, $\delta$ et $\alpha'$, $\beta'$, $\gamma'$,
$\delta'$ dans l'équation

\[ \text{III}\quad \alpha e^{\lambda\omega}-\beta e^{-\lambda\omega}+\gamma\cos.\lambda\omega-\delta\,\text{sin.}\,\lambda\omega=\alpha' e^{\lambda\omega}-\beta' e^{-\lambda\omega}+\gamma'\cos.\lambda\omega-\delta'\,\text{sin.}\,\lambda\omega \]

par cette substitution le premier membre devient:

\[ \frac{\alpha\{1-e^{\lambda\omega}(\text{sin.}\,\lambda\omega-\cos.\lambda\omega)\}+\alpha\{1+e^{-\lambda\omega}(\text{sin.}\,\lambda\omega+\cos.\lambda\omega)\}}{e^{-\lambda\omega}-\text{sin.}\,\lambda\omega+\cos.\lambda\omega} \]
\[ +\frac{\alpha\{e^{\lambda\omega}+e^{-\lambda\omega}+2\cos.\lambda\omega\}\cos.\lambda\omega+\alpha\{e^{\lambda\omega}-e^{-\lambda\omega}+2\,\text{sin.}\,\lambda\omega\}\,\text{sin.}\,\lambda\omega}{e^{-\lambda\omega}-\text{sin.}\,\lambda\omega+\cos.\lambda\omega} \]
\[ =\frac{\alpha\{\uncertain{4}-(e^{\lambda\omega}-e^{-\lambda\omega})\,\text{sin.}\,\lambda\omega+(e^{\lambda\omega}+e^{-\lambda\omega})\cos.\lambda\omega+(e^{\lambda\omega}+e^{-\lambda\omega})\cos.\lambda\omega+(e^{\lambda\omega}-e^{-\lambda\omega})\,\text{sin.}\,\lambda\omega\}}{e^{-\lambda\omega}-\text{sin.}\,\lambda\omega+\cos.\lambda\omega} \]
\[ =\frac{2\alpha\{2+(e^{\lambda\omega}+e^{-\lambda\omega})\cos.\lambda\omega\}}{e^{-\lambda\omega}-\text{sin.}\,\lambda\omega+\cos.\lambda\omega} \]

En faisant les substitutions dont il s'agit dans le second membre dela même
equation, il prend la forme.
\note{En tête, au milieu et souligné, le nombre 15 ; en haut à droite, à
l'encre, le nombre 122. Le premier terme entre accolades de l'avant-dernière
fraction est un chiffre repassé : la ligne suivante,
$2\alpha\{2+(e^{\lambda\omega}+e^{-\lambda\omega})\cos.\lambda\omega\}$,
demande 4, et c'est ce que les traits donnent le plus vraisemblablement ; la
lecture est signalée comme douteuse. Le feuillet, plus petit que la garde, est
monté sur un onglet ; il s'arrête sur « il prend la forme. », la suite étant à
la vue 252.}

\page{252}
\emph{16}

\[ \frac{\alpha' e^{(\lambda-1)\omega}\{e^{(1-\lambda)\omega}+\cos.(1-\lambda)\omega+\text{sin.}\,(1-\lambda)\omega\}+\alpha' e^{(1-\lambda)\omega}\{e^{(\lambda-1)\omega}+\cos.(1-\lambda)\omega-\text{sin.}\,(1-\lambda)\omega\}}{e^{-\omega}\{e^{(1-\lambda)\omega}+\cos.(1-\lambda)\omega+\text{sin.}\,(1-\lambda)\omega\}} \]
\[ +\frac{\alpha'\{(e^{(1-\lambda)\omega}-e^{(\lambda-1)\omega})\,\text{sin.}\,\omega+(e^{(1-\lambda)\omega}+e^{(\lambda-1)\omega})\cos.\omega+2\cos.\lambda\omega\}\cos.\lambda\omega}{e^{-\omega}\{e^{(1-\lambda)\omega}+\cos.(1-\lambda)\omega+\text{sin.}\,(1-\lambda)\omega\}} \]
\[ -\frac{\alpha'\{(e^{(1-\lambda)\omega}-e^{(\lambda-1)\omega})\cos.\omega-(e^{(1-\lambda)\omega}+e^{(\lambda-1)\omega})\,\text{sin.}\,\omega-2\,\text{sin.}\,\lambda\omega\}\,\text{sin.}\,\lambda\omega}{e^{-\omega}\{e^{(1-\lambda)\omega}+\cos.(1-\lambda)\omega+\text{sin.}\,(1-\lambda)\omega\}} \]
\[ =\frac{\alpha'\{\uncertain{4}+(e^{(1-\lambda)\omega}+e^{(\lambda-1)\omega})\cos.(1-\lambda)\omega-(e^{(1-\lambda)\omega}-e^{(\lambda-1)\omega})\,\text{sin.}\,(1-\lambda)\omega\}}{e^{-\omega}\{e^{(1-\lambda)\omega}+\cos.(1-\lambda)\omega+\text{sin.}\,(1-\lambda)\omega\}} \]
\[ +\frac{\alpha'\{(e^{(1-\lambda)\omega}+e^{(\lambda-1)\omega})[\cos.\lambda\omega\cos.\omega+\text{sin.}\,\lambda\omega\,\text{sin.}\,\omega]+(e^{(1-\lambda)\omega}-e^{(\lambda-1)\omega})[\text{sin.}\,\omega\cos.\lambda\omega-\cos.\omega\,\text{sin.}\,\lambda\omega]\}}{e^{-\omega}\{e^{(1-\lambda)\omega}+\cos.(1-\lambda)\omega+\text{sin.}\,(1-\lambda)\omega\}} \]
\[ =\frac{2\alpha'[2+(e^{(1-\lambda)\omega}+e^{(\lambda-1)\omega})\cos.(1-\lambda)\omega]}{e^{-\omega}\{e^{(1-\lambda)\omega}+\cos.(1-\lambda)\omega+\text{sin.}\,(1-\lambda)\omega\}} \]

Ainsi l'équation III se reduit en derniere analyse à

\[ \frac{\alpha\{2+(e^{\lambda\omega}+e^{-\lambda\omega})\cos.\lambda\omega\}}{e^{-\lambda\omega}-\text{sin.}\,\lambda\omega+\cos.\lambda\omega}=\frac{\alpha'\{2+(e^{(1-\lambda)\omega}+e^{(\lambda-1)\omega})\cos.(1-\lambda)\omega\}}{e^{-\omega}\{e^{(1-\lambda)\omega}+\cos.(1-\lambda)\omega+\text{sin.}\,(1-\lambda)\omega\}} \]

et par consequent
\[ \frac{\alpha e^{-\omega}\{e^{(1-\lambda)\omega}+\cos.(1-\lambda)\omega+\text{sin.}\,(1-\lambda)\omega\}}{\alpha'\{e^{-\lambda\omega}-\text{sin.}\,\lambda\omega+\cos.\lambda\omega\}}=\frac{2+(e^{(1-\lambda)\omega}+e^{(\lambda-1)\omega})\cos.(1-\lambda)\omega}{2+(e^{\lambda\omega}+e^{-\lambda\omega})\cos.\lambda\omega} \]

Il faut encore avoir recours a l'équation IV. \struck{que l'on peut mettre dans
laquelle} on peut faire passer tous les termes négatifs \add{de cette equation}
dans le second membre ce qui donne

\[ \alpha e^{\lambda\omega}+\beta e^{-\lambda\omega}-\gamma\,\text{sin.}\,\lambda\omega-\delta\cos.\lambda\omega=\alpha' e^{\lambda\omega}+\beta' e^{-\lambda\omega}-\gamma'\,\text{sin.}\,\lambda\omega-\delta'\cos.\lambda\omega \]

D'ou il resulte en ajoutant

\[ \text{I.}\quad \alpha e^{\lambda\omega}+\beta e^{-\lambda\omega}+\gamma\,\text{sin.}\,\lambda\omega+\delta\cos.\lambda\omega=0 \]
au premier membre et
\[ \text{II.}\quad \alpha' e^{\lambda\omega}+\beta' e^{-\lambda\omega}+\gamma'\,\text{sin.}\,\lambda\omega+\delta'\cos.\lambda\omega=0 \]
au second et divisant par 2.
\[ \alpha e^{\lambda\omega}+\beta e^{-\lambda\omega}=\alpha' e^{\lambda\omega}+\beta' e^{-\lambda\omega} \]

cette dernière équation devient en y mettant pour $\beta$ et $\beta'$ les
valeurs cidessus

\[ \frac{\alpha e^{\lambda\omega}\{e^{-\lambda\omega}-\text{sin.}\,\lambda\omega+\cos.\lambda\omega\}-\alpha e^{-\lambda\omega}\{e^{\lambda\omega}+\text{sin.}\,\lambda\omega+\cos.\lambda\omega\}}{e^{-\lambda\omega}-\text{sin.}\,\lambda\omega+\cos.\lambda\omega} \]
\[ =\frac{\alpha' e^{(\lambda-1)\omega}\{e^{(1-\lambda)\omega}+\cos.(1-\lambda)\omega+\text{sin.}\,(1-\lambda)\omega\}-\alpha' e^{(1-\lambda)\omega}\{e^{(\lambda-1)\omega}+\cos.(1-\lambda)\omega-\text{sin.}\,(1-\lambda)\omega\}}{e^{-\omega}\{e^{(1-\lambda)\omega}+\cos.(1-\lambda)\omega+\text{sin.}\,(1-\lambda)\omega\}} \]

ou
\[ \frac{\alpha\{(e^{\lambda\omega}-e^{-\lambda\omega})\cos.\lambda\omega-(e^{\lambda\omega}+e^{-\lambda\omega})\,\text{sin.}\,\lambda\omega\}}{e^{-\lambda\omega}-\text{sin.}\,\lambda\omega+\cos.\lambda\omega} \]
\[ =\frac{\alpha'\{(e^{(1-\lambda)\omega}+e^{(\lambda-1)\omega})\,\text{sin.}\,(1-\lambda)\omega-(e^{(1-\lambda)\omega}-e^{(\lambda-1)\omega})\cos.(1-\lambda)\omega\}}{e^{-\omega}\{e^{(1-\lambda)\omega}+\cos.(1-\lambda)\omega+\text{sin.}\,(1-\lambda)\omega\}} \]
\note{En tête, au milieu et souligné, le nombre 16 ; en haut à droite, à
l'encre, le nombre 123. Le premier terme entre accolades de la quatrième
fraction est le même chiffre repassé qu'à la vue 250 ; la réduction qui suit
demande 4, et la lecture est signalée comme douteuse. Le signe qui ouvre la
troisième fraction est un pâté épais, lu « $-$ ». Le feuillet, plus petit que
la garde, est monté sur un onglet ; il s'arrête sur cette dernière égalité,
qui se poursuit à la vue 254.}

\page{254}
\uncertain{\emph{17}}

D'ou on tire
\[ \frac{\alpha e^{-\omega}\{e^{(1-\lambda)\omega}+\cos.(1-\lambda)\omega+\text{sin.}\,(1-\lambda)\omega\}}{\alpha'\{e^{-\lambda\omega}-\text{sin.}\,\lambda\omega+\cos.\lambda\omega\}} \]
\[ =\frac{(e^{(1-\lambda)\omega}+e^{(\lambda-1)\omega})\,\text{sin.}\,(1-\lambda)\omega-(e^{(1-\lambda)\omega}-e^{(\lambda-1)\omega})\cos.(1-\lambda)\omega}{(e^{\lambda\omega}-e^{-\lambda\omega})\cos.\lambda\omega-(e^{\lambda\omega}+e^{-\lambda\omega})\,\text{sin.}\,\lambda\omega} \]

Il ne s'agit plus a present que de comparer ce resultat avec celui de
l'équation III pour conclure

\[ \{2+(e^{(1-\lambda)\omega}+e^{(\lambda-1)\omega})\cos.(1-\lambda)\omega\}\{(e^{\lambda\omega}+e^{-\lambda\omega})\,\text{sin.}\,\lambda\omega-(e^{\lambda\omega}-e^{-\lambda\omega})\cos.\lambda\omega\} \]
\[ +\{2+(e^{\lambda\omega}+e^{-\lambda\omega})\cos.\lambda\omega\}\{(e^{(1-\lambda)\omega}+e^{(\lambda-1)\omega})\,\text{sin.}\,(1-\lambda)\omega-(e^{(1-\lambda)\omega}-e^{(\lambda-1)\omega})\cos.(1-\lambda)\omega\}=0 \]

cette équation est celle dont il faut tirer la valeur de l'angle $\lambda\omega$.

Il est aisé de voir que dans le cas ou $\lambda=\frac{1}{2}$ elle se reduit a

\[ 2\{2+(e^{\frac{1}{2}\omega}+e^{-\frac{1}{2}\omega})\cos.\tfrac{1}{2}\omega\}\{(e^{\frac{1}{2}\omega}+e^{-\frac{1}{2}\omega})\,\text{sin.}\,\tfrac{1}{2}\omega-(e^{\frac{1}{2}\omega}-e^{-\frac{1}{2}\omega})\cos.\tfrac{1}{2}\omega\}=0 \]

et que cette équation est susceptible de deux solutions savoir.
\[ (e^{\frac{1}{2}\omega}+e^{-\frac{1}{2}\omega})\,\text{sin.}\,\tfrac{1}{2}\omega-(e^{\frac{1}{2}\omega}-e^{-\frac{1}{2}\omega})\cos.\tfrac{1}{2}\omega=0 \quad\text{et}\quad 2+(e^{\frac{1}{2}\omega}+e^{-\frac{1}{2}\omega})\cos.\tfrac{1}{2}\omega=0. \]

La premiere est relative au cas ou le stilet est \struck{simplement} consideré
comme simplement appuié au point du milieu dela lame dont les extremités sont
libres et elle montre que les deux moitiés de cette lame doivent vibrer comme
deux lames séparées de longueur $\frac{1}{2}a$ qui auroient chacune une de
leurs extremités libre et l'autre appuié, car la condition
$\dfrac{e^{\frac{1}{2}\omega}-e^{-\frac{1}{2}\omega}}{e^{\frac{1}{2}\omega}+e^{-\frac{1}{2}\omega}}=\text{tang.}\,\tfrac{1}{2}\omega$.
est celle qui convient au mouvement d'une \add{telle} lame
\struck{de longueur $\frac{1}{2}a$}

cette solution s'accorde au reste avec l'équation dependante de l'état des
extremités dela lame primitivee; car on a (V. cas. 1\textsuperscript{er})
$\cos.\omega=\dfrac{2e^{\omega}}{e^{2\omega}+1}$ D'ou on tire en prenant
$\text{sin.}\,\omega$ positif comme il convient ici,
$\text{sin.}\,\omega=\dfrac{e^{2\omega}-1}{e^{2\omega}+1}$ ainsi on a
$\cos.\omega=2\cos.^2\tfrac{1}{2}\omega-1=\dfrac{2e^{\omega}}{e^{2\omega}+1}$

\[ 2\cos.^2\tfrac{1}{2}\omega=\frac{(e^{\omega}+1)^2}{e^{2\omega}+1} \]
et
\[ 2\,\text{sin.}\,\tfrac{1}{2}\omega\cos.\tfrac{1}{2}\omega=\frac{e^{2\omega}-1}{e^{2\omega}+1} \]
\[ \text{tang.}\,\tfrac{1}{2}\omega=\frac{e^{\omega}-1}{e^{\omega}+1} \]

qui est precisement la condition resultante de la premiere solution.
\note{Le haut du feuillet est déchiré et taché : du nombre porté au milieu, en
tête, il ne subsiste que le bas des chiffres, lus avec doute « 17 » ; le
soulignement est entier. En haut à droite, à l'encre et intact, le nombre 124.
Dans la seconde fraction de la première égalité, le signe entre
$e^{\lambda\omega}$ et $e^{-\lambda\omega}$ et le mot « cos. » qui suit sont
empâtés. Après « le stilet est », un mot est biffé, lu « simplement », le même
mot revenant à la ligne suivante. « telle » est ajouté au-dessus de la ligne,
et « de longueur $\frac{1}{2}a$ » biffé après « lame ». « primitivee » est
écrit ainsi. La référence entre parenthèses est lue « V. cas.
1\textsuperscript{er} » ; la première lettre est tracée d'un seul jet et reste
douteuse. Le feuillet, plus petit que la garde, est monté sur un onglet ; il
s'arrête sur « solution. », la suite étant à la vue 256.}

\page{256}
\emph{18}

\struck{\ill{}}. A l'egard dela seconde \struck{équation} solution elle donne
$\dfrac{2}{e^{\frac{1}{2}\omega}+e^{-\frac{1}{2}\omega}}=-\cos.\tfrac{1}{2}\omega$
ainsi elle meneroit a conclure que la lame peut encore être consideré comme
séparé en deux autres lames de longueur $\frac{1}{2}a$ qui auroient chacune
une de leurs extremités fixée au point d'application du stilet et l'autre
libre, mais la condition
$\dfrac{2}{e^{\frac{1}{2}\omega}+e^{-\frac{1}{2}\omega}}=-\cos.\tfrac{1}{2}\omega$
ne peut s'accorder avec l'équation
$\cos.\omega=\dfrac{2e^{\omega}}{e^{2\omega}+1}$ dependante de l'état des
extremités que dans la seule supposition $\omega=0$. Il faut donc rejetter
cette seconde solution, comme il me semble necessaire dele faire dans le cas
examiné par Euler

\marginal{10}
Jusqu'ici, il existe la plus grande analogie entre le cas ou les deux
extremités sont appuiées et celui que je viens de développer où les deux
extremités sont libres car de part et d'autre la supposition
$\lambda=\frac{1}{2}$ donne des solutions semblables quant a la consideration
de \struck{point d'appui} \add{l'effet de l'application du stilet} mais
malheureusement un examen plus aprofondi du cas present fait apercevoir de
grandes differences.

En effet dans le cas des extremités appuiées la \struck{solu} premiere
solution relative a la supposition $\lambda=\frac{1}{2}$ n'est qu'un cas
particulier dela solution generale dans laquelle $\lambda$ peut
\struck{avoir toutes toutes les valeurs} \add{être une fraction quelconque
aliquote de l'unité} desorte que quelque soit celle adoptée la lame primitivee
peut toujours être consideré comme separé en deux lames des longueurs
$\lambda a$ et $(1-\lambda)a$ \struck{dont} qui auroient l'une et l'autre
leurs deux extremités appuiées, \struck{et} cette conclusion resulte de ce que
\struck{\ill{}} les equations $\text{sin.}\,\lambda\omega=0$ et
$\text{sin.}\,(1-\lambda)\omega=0$ ont lieu en même tems.
\struck{\ill{}} \add{pour de telles valeurs de $\lambda$ et de ce que} les
valeurs \struck{\ill{}} qui satisfont a ces equations sont les seules pour
lesquelles le mouvement dela lame soit possible.
\note{En tête, au milieu et souligné, le nombre 18 ; en haut à droite, à
l'encre, le nombre 125. La première ligne s'ouvre sur un mot ou un chiffre
biffé et empâté, suivi d'un point, qui n'est pas lu. Le chiffre 10 est seul
dans la marge de gauche, devant « Jusqu'ici » ; au-dessus, un petit signe
circulaire barré d'une croix. « l'effet de l'application » est ajouté au-dessus
de la ligne et « du stilet », qui l'achève, se lit dans la marge de gauche de
la ligne suivante. Les mots biffés avant « être une fraction quelconque
aliquote de l'unité » sont repassés les uns sur les autres ; la lecture
« avoir toutes toutes les valeurs » est donnée pour ce qu'elle vaut. Les trois
autres passages biffés des dernières lignes sont des pâtés et ne sont pas lus.
« primitivee » est écrit ainsi. Le feuillet, plus petit que la garde, est monté
sur un onglet ; il s'arrête sur « possible. », la suite étant à la vue 258.}

\page{258}
\emph{19}

\uncertain{Au} contraire dans le cas present des deux extremités libres les
équations
\[ \frac{e^{(1-\lambda)\omega}+e^{(\lambda-1)\omega}}{e^{(1-\lambda)\omega}-e^{(\lambda-1)\omega}}\,\text{sin.}\,(1-\lambda)\omega-\cos.(1-\lambda)\omega=0 \quad\text{et} \]
\[ \frac{e^{\lambda\omega}+e^{-\lambda\omega}}{e^{\lambda\omega}-e^{-\lambda\omega}}\,\text{sin.}\,\lambda\omega-\cos.\lambda\omega=0. \]
ne peuvent avoir lieu en même tems; lorsqu'on a égard a la condition qui
resulte de l'état des extremités de la lame primitivee, que dans le seul cas
ou $\lambda=\frac{1}{2}$ ces équations peuvent être mises sous la forme
\[ \text{tang.}\,(1-\lambda)\omega=\frac{e^{(1-\lambda)\omega}-e^{(\lambda-1)\omega}}{e^{(1-\lambda)\omega}+e^{(\lambda-1)\omega}}, \qquad \text{tang.}\,\lambda\omega=\frac{e^{\lambda\omega}-e^{-\lambda\omega}}{e^{\lambda\omega}+e^{-\lambda\omega}} \]
qui montre quelles ne sont autrechose que les conditions du mouvement de deux
lames des longueurs $\lambda a$ et $(1-\lambda)a$ dont une des extremités
seroit appuié et l'autre libre.

Ainsi \struck{ont} on est forcé de conclure que le cas ou le stilet est appuié
au milieu de la lame dont les deux extremités sont libres est le seul ou cette
lame puisse être consideré comme separé en deux autres lames qui auroient
chacune une de leurs extremités appuiees au point d'application du stilet et
l'autre libre.

\marginal{11}
Je vais maintenant démontrer qu'en effet les deux quantités
\[ \frac{e^{(1-\lambda)\omega}+e^{(\lambda-1)\omega}}{e^{(1-\lambda)\omega}-e^{(\lambda-1)\omega}}\,\text{sin.}\,(1-\lambda)\omega-\cos.(1-\lambda)\omega \quad\text{et}\quad \frac{e^{\lambda\omega}+e^{-\lambda\omega}}{e^{\lambda\omega}-e^{-\lambda\omega}}\,\text{sin.}\,\lambda\omega-\cos.\lambda\omega \]
ne peuvent être supposées en même tems égales a zéro.

c'est ceque montrent les transformations suivantes dans lesquelles on
introduit les valeurs tirées de la seconde dans la premiere des
\uncertain{avant} equations

\[ \frac{e^{(1-\lambda)\omega}+e^{(\lambda-1)\omega}}{e^{(1-\lambda)\omega}-e^{(\lambda-1)\omega}}\,\text{sin.}\,(1-\lambda)\omega-\cos.(1-\lambda)\omega \]
\[ =\frac{e^{(1-\lambda)\omega}+e^{(\lambda-1)\omega}}{e^{(1-\lambda)\omega}-e^{(\lambda-1)\omega}}\{\text{sin.}\,\omega\cos.\lambda\omega-\text{sin.}\,\lambda\omega\cos.\omega\}-\{\cos.\omega\cos.\lambda\omega+\text{sin.}\,\omega\,\text{sin.}\,\lambda\omega\} \]
\[ =\frac{e^{(1-\lambda)\omega}+e^{(\lambda-1)\omega}}{e^{(1-\lambda)\omega}-e^{(\lambda-1)\omega}}\left\{\frac{e^{\lambda\omega}+e^{-\lambda\omega}}{e^{\lambda\omega}-e^{-\lambda\omega}}\,\text{sin.}\,\lambda\omega\,\text{sin.}\,\omega-\text{sin.}\,\lambda\omega\cos.\omega\right\} \]
\[ -\left\{\frac{e^{\lambda\omega}+e^{-\lambda\omega}}{e^{\lambda\omega}-e^{-\lambda\omega}}\,\text{sin.}\,\lambda\omega\cos.\omega+\text{sin.}\,\omega\,\text{sin.}\,\lambda\omega\right\}=0 \]

et en multipliant par
$(e^{\lambda\omega}-e^{-\lambda\omega})(e^{(1-\lambda)\omega}-e^{(\lambda-1)\omega})$

\[ 0=\text{sin.}\,\lambda\omega\{(e^{\omega}+e^{(2\lambda-1)\omega}+e^{(1-2\lambda)\omega}+e^{-\omega})\,\text{sin.}\,\omega-(e^{\omega}+e^{(2\lambda-1)\omega}-e^{(1-2\lambda)\omega}-e^{-\omega})\cos.\omega \]
\[ -(e^{\omega}-e^{(2\lambda-1)\omega}+e^{(1-2\lambda)\omega}-e^{-\omega})\uncertain{\cos.\omega}-(e^{\omega}-e^{(2\lambda-1)\omega}-e^{(1-2\lambda)\omega}+e^{-\omega})\,\text{sin.}\,\omega\} \]

\struck{\ill{}}
\note{En tête, au milieu et souligné, le nombre 19 ; en haut à droite, à
l'encre, le nombre 126. Le chiffre 11 est seul dans la marge de gauche, devant
« Je vais maintenant ». Le premier mot du feuillet est tracé d'un jet : la
capitale se lit aussi bien O que A, et le sens demande « Au ». Le mot qui
précède « equations », après « des », n'est lu qu'avec doute : les jambages
donnent « avant ». Dans l'avant-dernière ligne, le « cos. $\omega$ » qui suit
la troisième parenthèse est couvert d'un pâté épais et n'est pas assuré.
La dernière ligne du feuillet est biffée d'un trait épais d'un bout à l'autre ;
on y distingue, sans certitude, « 2 sin. $\lambda\omega$
$(e^{\omega}-e^{-\omega})$ sin. $\omega$ ». « primitivee » est écrit ainsi.
Le feuillet, plus petit que la garde, est monté sur un onglet ; le calcul se
poursuit à la vue 260.}

\page{260}
\uncertain{\emph{20}}

\[ 0=2\,\text{sin.}\,\lambda\omega\{(e^{(1-2\lambda)\omega}+e^{(2\lambda-1)\omega})\,\text{sin.}\,\omega-(e^{\omega}-e^{-\omega})\cos.\omega\} \]

car amoins de faire $\lambda$ ou $\omega=0$, $\text{sin.}\,\lambda\omega$ ne
peut être supposé ici egal a zéro puisque cette \struck{hypothèse}
\add{supposition} meneroit a conclure que la portion de longueur $\lambda a$ a
ses deux extremités appuiée; contre l'hypothèse

Ainsi il faudroit que l'autre facteur
$(e^{(1-2\lambda)\omega}+e^{(2\lambda-1)\omega})\,\text{sin.}\,\omega-(e^{\omega}-e^{-\omega})\cos.\omega$
fut egal a zero D'ou on \uncertain{concluroit}
\[ \text{tang.}\,\omega=\frac{e^{\omega}-e^{-\omega}}{e^{(1-2\lambda)\omega}+e^{(2\lambda-1)\omega}} \]

mais la consideration de l'état des extremités dela lame primitivee donne
$\cos.\omega=\dfrac{2}{e^{\omega}+e^{-\omega}}$,
$\text{sin.}\,\omega=\pm\dfrac{e^{\omega}-e^{-\omega}}{e^{\omega}+e^{-\omega}}$
et par consequent
$\text{tang.}\,\omega=\pm\dfrac{e^{\omega}-e^{-\omega}}{2}$ on
\struck{seroit donc mené a ce resultat}

\marginal{*}
trouveroit donc en egalant les deux valeurs de $\text{tang.}\,\omega$
\[ \pm\frac{e^{\omega}-e^{-\omega}}{2}=\frac{e^{\omega}-e^{-\omega}}{e^{(1-2\lambda)\omega}+e^{(2\lambda-1)\omega}} \quad\text{ou}\quad \pm(e^{(1-2\lambda)\omega}+e^{(2\lambda-1)\omega})=2 \]

le signe superieur repond a la valeur positive de $\text{sin.}\,\omega$
c'est adire au cas ou il y a un nombre impair de nœud au contraire le signe
négatif repond a celui ou il y en a un nombre impair.

On voit d'abord que $e^{(1-2\lambda)\omega}+e^{(2\lambda-1)\omega}$ etant une
quantité essentiellement positive l'équation
$-(e^{(1-2\lambda)\omega}+e^{(2\lambda-1)\omega})=2$ ne peut subsister dans
aucun cas et a l'egard de celle ci
$e^{(1-2\lambda)\omega}+e^{(2\lambda-1)\omega}=2$ on voit qu'amoins de faire
$\omega=0$, ceque exclut tout espece de mouvement, elle ne peut être
satisfaite pour aucune valeur de $\lambda$ differente de $\frac{1}{2}$

\marginal{12}
On obtient des resultats semblables pour le cas ou les extremités de la lame
aulieu d'être libres sont fixées car si on suppose qu'une telle lame soit
soumise a l'action d'un stilet dans un point quelconque de son etendue on
parvient, au moyen de calculs que je suprime ici a cause de leur parfaite
analogie avec les précédans, a l'équation.
\note{Deux nombres soulignés en tête : sur le feuillet lui-même, au milieu, un
nombre dont le second chiffre est tracé comme un $\omega$ et qui est lu, avec
doute, 20 ; au-dessus du bord supérieur du feuillet, donc sur le montage et
non sur cette page, un autre nombre souligné qui se lit 21. En haut à droite,
à l'encre et sur le feuillet, le nombre 127. « supposition » est ajouté
au-dessus de « hypothèse », biffé. Le « roit » de « concluroit » est couvert
d'un pâté. Devant « trouveroit », dans la marge, un petit signe en étoile qui
rattache la phrase à « on », resté au bout de la ligne précédente. La phrase
« au contraire le signe négatif repond a celui ou il y en a un nombre impair »
répète « impair » : le leçon est celle du feuillet, et le raisonnement
demanderait « pair ». Le feuillet, plus petit que la garde, est monté sur un
onglet ; il s'arrête sur « a l'équation. », et la suite est au-delà de ce lot.}

\end{document}
